\documentclass[11pt,a4paper]{article}

% ----------------------------------------------------------------
%  Paquetes Avanzados
% ----------------------------------------------------------------
\usepackage[utf8]{inputenc}
\usepackage[T1]{fontenc}
\usepackage[spanish,es-nodecimaldot]{babel}
\usepackage[margin=2.5cm]{geometry}
\usepackage{amsmath,amssymb,amsfonts}
\usepackage{booktabs}
\usepackage{xcolor}
\usepackage{graphicx}
\usepackage{enumitem}
\usepackage{listings}
\usepackage{tcolorbox}
\usepackage{tikz}
\usetikzlibrary{shapes.geometric, arrows.meta, positioning, calc, fit, backgrounds, shadows.blur}
\usepackage[hidelinks]{hyperref}
\usepackage{xurl}

\definecolor{codegray}{gray}{0.95}
\definecolor{primary}{RGB}{41, 128, 185}
\definecolor{danger}{RGB}{192, 57, 43}
\definecolor{success}{RGB}{39, 174, 96}
\definecolor{warning}{RGB}{243, 156, 18}

\lstset{
  basicstyle=\ttfamily\footnotesize,
  backgroundcolor=\color{codegray},
  breaklines=true,
  frame=single,
  columns=fullflexible,
  showstringspaces=false,
}

\tcbuselibrary{skins,breakable}
\newtcolorbox{pregunta}[2][]{
    colback=primary!5!white,
    colframe=primary!80!black,
    fonttitle=\bfseries,
    title={Pregunta: #2},
    breakable,
    drop shadow,
    #1
}

\newtcolorbox{trampa}[2][]{
    colback=danger!5!white,
    colframe=danger!80!black,
    fonttitle=\bfseries,
    title={Pregunta Trampa: #2},
    breakable,
    drop shadow,
    #1
}

\title{\textbf{\Huge Wiki de Conocimiento y Defensa}\\
\vspace{0.5cm}
\Large Detecci\'on (YOLOv11) + OCR CNN + L\'ogica Difusa Takagi-Sugeno}
\author{\textbf{Material de Estudio para Defensa del Proyecto} \\ Inteligencia Artificial --- Universidad T\'ecnica de Ambato}
\date{Junio 2026}

\begin{document}
\maketitle
\tableofcontents
\newpage

% ================================================================
\section{Arquitectura del Sistema y Decisiones de Dise\~no}
% ================================================================

\subsection{Pipeline de Procesamiento (Diagrama)}

El sistema emplea una arquitectura \textbf{h\'ibrida}. En lugar de usar una red neuronal \emph{end-to-end} (que recibe la imagen y arroja el texto), el problema se divide en etapas especializadas. Esto garantiza mayor interpretabilidad, facilidad de depuraci\'on y determinismo.

\begin{figure}[h!]
    \centering
    \begin{tikzpicture}[
        node distance=1.5cm and 2.5cm,
        box/.style={rectangle, draw, rounded corners, fill=white, drop shadow, align=center, minimum height=1cm, minimum width=2.5cm, font=\sffamily\footnotesize},
        arrow/.style={-{Stealth[scale=1.2]}, thick},
        group/.style={rectangle, draw, dashed, inner sep=0.3cm}
    ]
        % Nodes
        \node[box, fill=gray!10] (cam) {Fotograma\\(C\'amara)};
        \node[box, right=of cam, fill=blue!10] (yolo) {YOLOv11\\(Detecci\'on de Placa)};
        \node[box, right=of yolo, fill=orange!10] (prep) {Preproceso\\(CLAHE + Gris)};
        \node[box, right=of prep, fill=orange!10] (seg) {Segmentaci\'on\\(Umbral Gaussiano)};
        \node[box, below=of seg, fill=purple!10] (cnn) {CNN Clasificadora\\(36 clases, Batch)};
        \node[box, left=of cnn, fill=green!10] (decod) {Decodificaci\'on\\Posicional};
        \node[box, left=of decod, fill=green!10] (vot) {Votaci\'on Temporal\\(Multi-frame)};
        \node[box, below=of vot, fill=red!10] (difuso) {L\'ogica Difusa\\(Takagi-Sugeno)};
        \node[box, right=of difuso, fill=teal!10] (json) {Registro JSON\\+ Email};

        % Backgrounds for modules
        \begin{scope}[on background layer]
            \node[group, fit=(yolo) (prep) (seg) (cnn), fill=yellow!5, label={[anchor=south]above:\textbf{M\'odulo de Visi\'on (OCR S\'incrono)}}] {};
            \node[group, fit=(decod) (vot), fill=green!5, label={[anchor=south]above:\textbf{Post-procesamiento}}] {};
        \end{scope}

        % Arrows
        \draw[arrow] (cam) -- (yolo);
        \draw[arrow] (yolo) -- node[above, font=\tiny] {Recorte Bounding Box} (prep);
        \draw[arrow] (prep) -- (seg);
        \draw[arrow] (seg) -- node[right, font=\tiny] {Cajas de Caracteres} (cnn);
        \draw[arrow] (cnn) -- node[above, font=\tiny] {Softmax Probs} (decod);
        \draw[arrow] (decod) -- (vot);
        \draw[arrow] (vot) -- node[left, font=\tiny] {Placa Final + Velocidad} (difuso);
        \draw[arrow] (difuso) -- node[above, font=\tiny] {Sanci\'on (Horas)} (json);

    \end{tikzpicture}
    \caption{Arquitectura del Sistema de Radar y OCR.}
\end{figure}

\subsection{El Debate Cr\'itico: OCR S\'incrono vs As\'incrono}
Una de las decisiones arquitect\'onicas m\'as fuertes fue \textbf{mantener el OCR en modo s\'incrono} (procesado en el bucle principal) y no enviar la tarea a un \texttt{ThreadPoolExecutor} as\'incrono.

\begin{itemize}
    \item \textbf{El problema as\'incrono:} En un sistema con una sola GPU, el tracker de objetos (YOLO con ByteTrack) y el clasificador CNN compiten por los recursos CUDA. Si se usan hilos (as\'incrono), la alternancia de contexto en la GPU genera cuellos de botella no deterministas (latencia espasm\'odica) y p\'erdida de fotogramas (lag).
    \item \textbf{La soluci\'on s\'incrona:} Se procesa exactamente \textbf{1 OCR por frame}. Para no desperdiciar recursos, las detecciones de YOLO se \emph{ordenan por \'area}. As\'i, el "presupuesto" de 1 OCR se gasta siempre en el veh\'iculo \textbf{m\'as cercano y grande}, garantizando la m\'axima calidad de resoluci\'on de la placa sin bloquear la GPU para el tracking del resto de autos.
\end{itemize}

% ================================================================
\section{Visi\'on por Computador Avanzada}
% ================================================================

\subsection{Preprocesamiento y Resoluciones: El Umbral de 48px}
El sistema extrae el \emph{bounding box} de la placa directamente del fotograma original en alta resoluci\'on, no del tensor reescalado ($416\times416$) de YOLO.

Para el OCR, el tama\~no importa. Si la placa mide menos de \textbf{48 p\'ixeles de altura}, se aplica un \textbf{escalado bic\'ubico} (\texttt{cv2.INTER\_CUBIC}).
\begin{itemize}
    \item \textbf{¿Por qu\'e 48px?} Se baj\'o el umbral hist\'orico de 72px a 48px emp\'iricamente. Escalar placas medianas (50-70px) a\~nad\'ia ruido de interpolaci\'on innecesario.
    \item \textbf{¿Por qu\'e Bic\'ubico y no Vecino M\'as Pr\'oximo?} El vecino m\'as pr\'oximo genera artefactos de "escalera" (aliasing) que destruyen las curvas de las letras (O, C, 8). La bic\'ubica usa una f\'ormula polinomial sobre una vecindad $4\times4$, preservando gradientes suaves en los bordes para la posterior binarizaci\'on.
\end{itemize}

\subsection{Binarizaci\'on: Umbral Adaptativo Gaussiano vs Otsu}
\textbf{Otsu Global} calcula un umbral \'unico minimizando la varianza intra-clase para toda la imagen. Falla estrepitosamente en placas vehiculares porque la iluminaci\'on nunca es uniforme (sombras de parachoques, reflejos del sol). Otsu terminaba oscureciendo media placa y fusionando caracteres.

El sistema usa \textbf{Umbral Adaptativo Gaussiano}. El umbral se calcula din\'amicamente para un p\'ixel bas\'andose en el promedio ponderado (gaussiano) de su vecindad ($N \times N$). Si el bloque est\'a en sombra, el umbral local baja; si est\'a al sol, sube. Esto independiza los caracteres del gradiente de luz del fondo.

\subsection{Decodificaci\'on Posicional Categ\'orica}
La red CNN arroja probabilidades \emph{Softmax} sobre 36 clases ($0-9$, $A-Z$). Existen letras y n\'umeros visualmente indiscernibles a baja resoluci\'on ($O \leftrightarrow 0$, $I \leftrightarrow 1$, $Z \leftrightarrow 2$, $B \leftrightarrow 8$).

Dado que la placa ecuatoriana tiene el patr\'on estricto \texttt{ABC-NNNN} (3 letras, 3-4 d\'igitos), aplicamos un \textbf{argmax segmentado}:
\[
c_i =
\begin{cases}
\arg\max_{k\in\mathcal{L}} p^{(i)}_k & \text{si } i < 3 \quad (\mathcal{L} = \text{Letras})\\
\arg\max_{k\in\mathcal{D}} p^{(i)}_k & \text{si } i \geq 3 \quad (\mathcal{D} = \text{D\'igitos})
\end{cases}
\]
Si la CNN arroja $p(0)=0.8, p(O)=0.2$ en la posici\'on $i=1$, el sistema fuerza la decodificaci\'on a $O$, corrigiendo el error intr\'inseco del modelo.

% ================================================================
\section{L\'ogica Difusa: Las Matem\'aticas de la Sanci\'on}
% ================================================================

Este proyecto implementa una arquitectura difusa \emph{h\'ibrida} (Mamdani para categorizaci\'on y Takagi-Sugeno para sanci\'on continua). Es crucial entender por qu\'e se abandon\'o un sistema Mamdani puro para la sanci\'on en horas.

\subsection{El problema del "Serrucho" (P\'erdida de Monotonicidad en Mamdani)}
En un inicio, la salida "horas de indisponibilidad" se model\'o con Mamdani (defuzzificaci\'on por centroide de \'areas).
Si tenemos reglas consecutivas (Multa Leve $\rightarrow$ baja, Multa Moderada $\rightarrow$ media), y calculamos el \textbf{centroide de la uni\'on de \'areas recortadas}, ocurren ca\'idas abruptas locales (el efecto "serrucho").
Cuando la velocidad aumenta y empieza a activar la siguiente regla, el \'area de la nueva regla al principio es peque\~na, "jalando" el centroide hacia la zona compartida, lo que matem\'aticamente produce que a \textbf{mayor velocidad, haya ligeramente menor sanci\'on localmente}. Esto viola el principio legal de proporcionalidad.

\subsection{La Soluci\'on: Takagi-Sugeno de Orden Cero}
El modelo de Sugeno descarta las \'areas de salida. El consecuente de la regla $i$ es un valor constante (un \emph{singleton}) $c_i$. La salida total $H(v)$ es el promedio de los singletons ponderado por el grado de activaci\'on $\mu_i(v)$ de cada regla:

\[
H(v) = \frac{\sum_{i=1}^{n} \mu_i(v) \cdot c_i}{\sum_{i=1}^{n} \mu_i(v)}
\]

Dado que los singletons son estrictamente crecientes ($c_1 < c_2 < \dots < c_n$) y los antecedentes teselan el espacio (tri\'angulos solapados 50\%), la curva resultante $H(v)$ es \textbf{estrictamente mon\'otona y suave}. A m\'as velocidad, \textbf{siempre} hay m\'as sanci\'on.

\subsection{Las 5 Reglas Expl\'icitas del Sistema de Sanci\'on}
Solo cuando el sistema Mamdani de nivel superior detecta la categor\'ia \textbf{Multa} ($>20$ km/h), se aplica la evaluaci\'on Sugeno de horas:

\begin{lstlisting}[language=Python]
_REGLAS_SUGENO = [
    # (nombre,        [params MF_entrada],   singleton_salida_horas)
    ("multa_leve",     [20, 24, 28],            5),   # Exceso  1-8  km/h -> 5h
    ("multa_moderada", [24, 28, 32],           18),   # Exceso  5-12 km/h -> 18h
    ("multa_alta",     [28, 32, 36],           34),   # Exceso  9-16 km/h -> 34h
    ("multa_grave",    [32, 36, 40],           52),   # Exceso 13-20 km/h -> 52h
    ("multa_extrema",  [36, 40, 40, 40],       70),   # Exceso >=16  km/h -> 70h
]
\end{lstlisting}

% ================================================================
\section{Glosario R\'apido de Conceptos de IA}
% ================================================================
\begin{description}[leftmargin=!,labelwidth=4cm]
    \item[\textbf{Epoch (Época)}] Una pasada completa de \emph{todos} los datos de entrenamiento a trav\'es de la red neuronal.
    \item[\textbf{Batch Size}] El n\'umero de ejemplos procesados antes de actualizar los pesos del modelo. Usamos 256 para estabilidad del gradiente.
    \item[\textbf{Learning Rate (Tasa de aprendizaje)}] El tama\~no del paso de correcci\'on de los pesos. Se usa un \emph{planificador OneCycleLR} (empieza bajo, sube a un pico, y baja asint\'oticamente).
    \item[\textbf{AMP (Automatic Mixed Precision)}] Uso de formato FP16 (16 bits) en vez de FP32 (32 bits) en la GPU. Duplica la velocidad de entrenamiento y reduce el uso de VRAM a la mitad sin perder precisi\'on algor\'itmica.
    \item[\textbf{Label Smoothing}] En vez de forzar a la red a predecir 100\% de confianza para la respuesta correcta (vector $[0, 0, 1, 0]$), se asigna un 95\% a la correcta y 5\% distribuido al resto. Esto evita sobreconfianza y \emph{overfitting}.
\end{description}

\newpage
% ================================================================
\section{Banco de Preguntas para la Defensa (Simulador)}
% ================================================================
Esta secci\'on contiene el n\'ucleo para la defensa. El profesor intentar\'a buscar huecos en la l\'ogica del estudiante.

\begin{trampa}{¿Por qu\'e no entrenaron una red (CNN/YOLO) para estimar la velocidad directamente de la imagen?}
    \textbf{Respuesta:} "Porque estimar velocidad desde un sensor monocular (1 sola c\'amara) mediante deep learning puro es un problema sub-determinado y prono a errores masivos."\\
    Un modelo CNN requerir\'ia mapeo 3D a 2D (estimaci\'on de profundidad y matriz de calibraci\'on intr\'inseca de la c\'amara) que var\'ia dependiendo de cu\'anto zoom se aplique o el tama\~no del auto. Nuestro enfoque geom\'etrico con dos \textbf{l\'ineas virtuales de distancia conocida} se rige por F\'isica Cl\'asica ($v = \Delta d / \Delta t$) y garantiza determinismo absoluto sin requerir calibraci\'on est\'ereo ni lidar.
\end{trampa}

\begin{pregunta}{¿Por qu\'e desecharon Otsu en la segmentaci\'on de caracteres?}
    \textbf{Respuesta:} "El m\'etodo Otsu asume que la imagen tiene un histograma bimodal estricto (fondo y texto bien separados). En un ambiente al aire libre, un extremo de la placa puede tener sol directo y el otro la sombra del maletero del auto. El umbral global de Otsu promedia esto y destruye el texto en la sombra o quema el texto en el sol. Usamos \textbf{Umbral Adaptativo Gaussiano} porque calcula el umbral p\'ixel a p\'ixel en funci\'on de su vecindad local ($21\times 21$ píxeles), mitigando por completo las gradientes de luz."
\end{pregunta}

\begin{trampa}{Si su l\'ogica difusa es tan precisa para dar las horas de sanci\'on, ¿por qu\'e no usaron una funci\'on matem\'atica simple, tipo $f(x) = (x - 20) \times 3$?}
    \textbf{Respuesta:} "Por la \textbf{capacidad de modelado experto} que ofrece la l\'ogica difusa." \\
    Una funci\'on lineal o polinomial trata todas las infracciones con la misma sensibilidad. La l\'ogica difusa nos permite modelar la \textbf{heur\'istica humana y legal}: la transici\'on de leve a moderado requiere un nivel de severidad distinto al de grave a extremo. Al definir conjuntos difusos traslapados, podemos establecer mesetas, rampas de transici\'on suave y umbrales no-lineales basados en reglas l\'ingü\'isticas (IF 'velocidad es alta' THEN 'sanción es moderada'), lo cual es imposible de codificar de forma interpretable con una simple f\'ormula $f(x)$ o con r\'igidos If-Else que generan saltos abruptos.
\end{trampa}

\begin{pregunta}{En su clasificador CNN usaron Softmax. ¿Cu\'al es la diferencia entre Softmax y Sigmoide y por qu\'e no usaron Sigmoide?}
    \textbf{Respuesta:} "La Sigmoide comprime cualquier n\'umero al rango $[0,1]$ de forma independiente, asumiendo que un input podr\'ia pertenecer a m\'ultiples clases a la vez (\emph{Multi-label classification}). El \textbf{Softmax} escala las salidas de modo que \textbf{la suma de todas las probabilidades de las 36 clases sea exactamente 1.0} (\emph{Multi-class classification}). Como un car\'acter en la placa solo puede ser \textbf{una sola cosa} (es una 'A' o es una 'B', no ambas), Softmax obliga a las probabilidades a competir entre s\'i, dando una distribuci\'on de confianza real."
\end{pregunta}

\begin{trampa}{Si imprimo una placa falsa en un papel y camino frente a su c\'amara, ¿me multan?}
    \textbf{Respuesta:} "S\'i y no. YOLO detectar\'ia la placa porque su patr\'on visual existe (bordes, texto). Sin embargo, el \textbf{filtro de clases} de nuestro tracker asocia la detecci\'on al \emph{objeto padre} (Auto, Moto, Camión, Bus en COCO). Si una persona lleva el cart\'on, YOLO clasifica a la persona como \texttt{person} (clase 0) y no como veh\'iculo (clases 2,3,5,7). Como solo medimos cruce de l\'ineas para veh\'iculos vivos, el peat\'on con el cart\'on no disparar\'ia el flujo de velocidad."
\end{trampa}

\begin{pregunta}{He visto que su modelo CNN fue entrenado con placas de la India. ¿C\'omo es posible que funcione en placas Ecuatorianas?}
    \textbf{Respuesta:} "Esta es una distinci\'on clave de nuestra arquitectura. Si hubi\'eramos usado una red OCR \emph{End-to-End} (que lee toda la imagen de corrido, como CRNN), el modelo habr\'ia aprendido a leer el formato y la fuente de la India, fallando en Ecuador. 
    Nosotros dividimos el problema: primero extraemos los caracteres sueltos mediante algoritmos de visi\'on (Contornos y Componentes Conexos) y nuestra CNN s\'olo clasifica recortes individuales de $48\times 48$. \textbf{El n\'umero '5' o la letra 'M' de una placa de la India es tipogr\'aficamente id\'entico al de Ecuador}. El modelo de \emph{Deep Learning} solo reconoce los glifos universales; la l\'ogica del formato ecuatoriano (ABC-NNNN) la aplicamos nosotros en el software."
\end{pregunta}

\begin{trampa}{Me dicen que usan "Votaci\'on Temporal" porque el OCR falla. ¿Su red neuronal es mala?}
    \textbf{Respuesta:} "No, nuestra red tiene un \textbf{93.6\% de precisi\'on por car\'acter}, lo cual es excelente considerando el ruido exterior. El problema es estoc\'astico: si una placa tiene 7 caracteres, la probabilidad de leer la placa perfecta en un s\'olo \emph{frame} es $0.936^7 \approx 62\%$. 
    Sin embargo, el veh\'iculo aparece en nuestro campo de visi\'on durante unos 10-15 frames. El error causado por el ruido (\emph{motion blur}, sol, una gota de lluvia) es \textbf{aleatorio}, mientras que el car\'acter correcto es \textbf{constante}. Al acumular las lecturas y aplicar la Moda Matem\'atica por columna (Votaci\'on Temporal), el ruido aleatorio se cancela y la se\~nal constante emerge, d\'andonos secuencias correctas con una precisi\'on pr\'actica cercana al 100\%."
\end{trampa}

\begin{pregunta}{Veo que usan $lr=10^{-3}$ y AdamW. ¿Por qu\'e AdamW en vez del m\'as famoso Adam regular?}
    \textbf{Respuesta:} "AdamW (Adam con Weight Decay acoplado) soluciona un error matem\'atico en la implementaci\'on original de Adam en PyTorch. Adam aplica la regularizaci\'on L2 (\emph{weight decay}) dentro del c\'alculo de los promedios m\'oviles del gradiente, lo que hace que la regularizaci\'on pierda fuerza, causando sobreajuste. \textbf{AdamW separa el \emph{weight decay}} sum\'andolo directamente a los pesos despu\'es de la optimizaci\'on de gradiente, logrando que el modelo generalice mejor frente a datos no vistos."
\end{pregunta}

\begin{trampa}{En su \texttt{pipeline.py} vi que asignan un 'presupuesto' de OCR = 1 por frame. ¿Qu\'e pasa si pasan dos veh\'iculos exactamente al mismo tiempo y cruzan la l\'inea?}
    \textbf{Respuesta:} "Ambos ser\'an detectados y su velocidad ser\'a medida, ya que YOLO y el Tracker (ByteTrack) son algoritmos paralelos y calculan el \emph{bounding box} de todos en cada frame. 
    Respecto a la placa, el presupuesto de $OCR=1$ hace que el sistema clasifique la placa del auto que tenga el \textbf{\'area m\'as grande en pantalla} (sort por \'area). Como los autos cruzan en milisegundos distintos, el auto m\'as cercano consume el OCR durante unos frames, y una vez que pasa, el siguiente auto se vuelve el de mayor \'area y hereda el presupuesto. Como tenemos de 10 a 15 frames efectivos en la zona de c\'amara lenta por segundo, hay frames m\'as que suficientes para atender a ambos de forma s\'incrona sin bloquear el sistema."
\end{trampa}

\begin{pregunta}{¿Qu\'e es el hiperpar\'ametro 'IoU' en YOLO y por qu\'e es crucial para su tracker?}
    \textbf{Respuesta:} "IoU (Intersection over Union) es la medida de solapamiento entre dos Bounding Boxes (Área de Intersecci\'on dividida para el Área de Uni\'on). El tracker ByteTrack lo usa para el \emph{Data Association}: si en el frame $T$ hay una caja de auto y en el frame $T+1$ hay otra, el algoritmo calcula su IoU espacial. Si es muy alto (las cajas se solapan), el tracker concluye que es exactamente \textbf{el mismo auto} movi\'endose y le mantiene su ID único. Sin IoU, el tracker no sabr\'ia c\'omo unir las detecciones entre \emph{frames}."
\end{pregunta}

\end{document}
